%TCIDATA{Version=5.00.0.2606}
%TCIDATA{LaTeXparent=0,0,main-sw.tex}
                      
%TCIDATA{ChildDefaults=chapter:1,page:1}


\section{Instructions for Using the Thesis Shells}

We give a brief list of instructions for using the thesis shells with
Scientific Workplace, some of which are listed on the department's web page.

You will need to download all of the front matter files along with
nmsuth01.cls, main-sw.tex, and thesisbody.tex. Place all of these files in
the same directory. You should also read the file main-sw.dvi or main-sw.ps
or main-sw.pdf for more detailed information about using these files.

Write the body of your thesis, using any style you wish. You may also write
an abstract page in Scientific Workplace. Once you are finished with your
typing, open the file thesisbody.tex, delete any text you find, except for
the grey boxes right before the sample references. Then click on file, then
import contents, and then locate the file for your thesis. Then save the
file. Next, edit the files title-sw.tex, approval-sw.tex, dedica-sw.tex,
ackno-sw.tex, vita-sw.tex, and abstract-sw.tex, by putting in the
appropriate data. Make sure you do not delete any of the blank spaces or
other formatting commands. If you typed an abstract, open the file
abstract-sw.tex, place the cursor immediately before the first line of the
abstract itself (which starts Type your abstract here), click on file, then
import contents, and enter the name of your abstract file. Finally, delete
the existing line starting with Type your abstract here, and then save.

To compile your document, click on typeset, then preview, and change the
number of passes to three. You should only compile main-sw, not thesisbody.

\subsection{Notes for Typing}

\begin{enumerate}
\item Do not use chapter headings in your thesis. The most major heading
should be a section heading. You will need to change your headings if you
used chapter heads.

\item Be careful when editing the front matter documents; the formatting
commands need to be kept as they are if you want your thesis correctly
formatted.

\item If you do not have tables and/or figures, delete the appropriate line
in the main-sw.tex file.

\item If you typed your thesis in several parts, import all of them into the
file thesisbody.tex.

\item You are encouraged to read the file main.dvi or main.ps or main.pdf
for more detailed information about typing your thesis.
\end{enumerate}

\subsection{What do you do if you have a really long heading, such as this
long subsection heading?}

If you have headings that are more than one line in the table of contents,
then thesisbody.tex will need to be modified. Perform the following steps:

\begin{enumerate}
\item preview your thesis to find out which headings take more than one
line. For each, determine where the first line of the heading ends.

\item at the location where the first line of a long heading ends, add a tex
field (click on insert, then field, then tex field) with the following
command: \texttt{\TEXTsymbol{\backslash}vspace\{-\TEXTsymbol{\backslash}%
singlespace\} \TEXTsymbol{\backslash}\TEXTsymbol{\backslash}}

\item before clicking ok, click the box named Encapsulated at the top left
of the tex field screen.
\end{enumerate}

To see how this works, when you preview this document, you will see the
first line of the section head ending after the work this. We then modify
the section heading as follows.

\subsection{What do you do if you have a really long heading, such as this%
%TCIMACRO{\TeXButton{spacing inside heading}{\vspace{-\singlespace} \\ } }%
%BeginExpansion
\vspace{-\singlespace} \\ 
%EndExpansion
long subsection heading?}

You will see, in the table of contents and the document itself the
difference between the two versions of the subsection heads.

\subsection{Figures and Tables}

If you have figures and or tables in your document, and you want them to
show up in a list of tables and list of figures, there are a few things you
must do. First, any graphic that is to show up in the list of figures should
be a floating graphic. If you have a graphic, click on it to open the small
blue icon at the bottom right of the graphic, then click on the icon. Click
on the layout tab, and click the button labeled floating. For tables, make
them by adding the fragment named \emph{Table - (4 x 3, floating)}. Open the
grey box named caption and type in your caption, making sure not to change
any of the tex code. If you want to refer to the table, change the text
inside the \TEXTsymbol{\backslash}label command. If you have already typed a
table, add the fragment, delete the sample table, and place the grey boxes
around your table in the same position as in the sample.

If you do not have figures and or tables, you should edit the main-sw.tex
file. You will see the following grey boxes:

\begin{center}
%TCIMACRO{\TeXButton{newpage}{ }}%
%BeginExpansion
 %
%EndExpansion
%TCIMACRO{\TeXButton{List of Tables}{ }}%
%BeginExpansion
 %
%EndExpansion
%TCIMACRO{\TeXButton{addcontentsline List of Tables}{ }}%
%BeginExpansion
 %
%EndExpansion

%TCIMACRO{\TeXButton{newpage}{}}%
%BeginExpansion
%
%EndExpansion
%TCIMACRO{\TeXButton{List of Figures}{ }}%
%BeginExpansion
 %
%EndExpansion
%TCIMACRO{\TeXButton{addcontentsline List of Figures}{ }}%
%BeginExpansion
 %
%EndExpansion
\end{center}

Delete the appropriate entire line of these boxes if you do not have tables
or figures.

Here is an example of a figure and a table.\FRAME{ftbpFU}{3.1955in}{2.3774in%
}{0pt}{\Qcb{example figure}}{}{sampleplot.jpg}{\special{language "Scientific
Word";type "GRAPHIC";maintain-aspect-ratio TRUE;display "USEDEF";valid_file
"F";width 3.1955in;height 2.3774in;depth 0pt;original-width
4.0413in;original-height 3in;cropleft "0";croptop "1";cropright
"1";cropbottom "0";filename 'SamplePlot.jpg';file-properties "XNPEU";}}

\begin{center}
%TCIMACRO{\TeXButton{B}{\begin{table}[tbp] \centering}}%
%BeginExpansion
\begin{table}[tbp] \centering%
%EndExpansion
\begin{tabular}{|l|l|l|}
\hline
\textbf{Head} & \textbf{Head} & \textbf{Head} \\ \hline
entry & entry & entry \\ 
entry & entry & entry \\ 
entry & entry & entry \\ \hline
\end{tabular}%
\caption{example table\label{key}}%
%TCIMACRO{\TeXButton{E}{\end{table}}}%
%BeginExpansion
\end{table}%
%EndExpansion
\end{center}


\begin{center}
%TCIMACRO{\TeXButton{B}{\begin{table}[tbp] \centering}}%
%BeginExpansion
\begin{table}[tbp] \centering%
%EndExpansion
\begin{tabular}{|l|l|l|}
\hline
\textbf{Head} & \textbf{Head} & \textbf{Head} \\ \hline
entry & entry & entry \\ 
entry & entry & entry \\ 
entry & entry & entry \\ \hline
\end{tabular}%
\caption{example table two\label{key}}%
%TCIMACRO{\TeXButton{E}{\end{table}}}%
%BeginExpansion
\end{table}%
%EndExpansion
\end{center}

\subsection{Footnotes}

If you have long footnotes\footnote{%
%TCIMACRO{\TeXButton{singlespace}{\setlength{\baselineskip}{\singlespace}}}%
%BeginExpansion
\setlength{\baselineskip}{\singlespace}%
%EndExpansion
What do you do if you have a really long footnote and you do not wish to
make it smaller? The graduate school will want footnotes single spaced, and
if you have more than one footnote on a page, doublespacing between
footnotes.%
%TCIMACRO{\TeXButton{vspace}{\vspace{\singlespace}}}%
%BeginExpansion
\vspace{\singlespace}%
%EndExpansion
}, you will have to reformat them since the graduate school wants footnotes
single spaced. Also, if you have more than one footnote on a page, they want
doublespacing between footnotes. To accomplish this, at the beginning of the
footnote box, insert a tex field (click on insert, typeset object, tex
field, or some similar set of menu items if you are using a version of
Scientific Workplace earlier than 4.0). In the tex field, add the line

\begin{center}
\TEXTsymbol{\backslash}setlength\{\TEXTsymbol{\backslash}baselineskip\}\{%
\TEXTsymbol{\backslash}singlespace\}
\end{center}

\noindent You must do this to each footnote. If you have more than one
footnote on a page, to each but the last footnote on the page add a tex
field in the footnote box after the footnote text containing the line

\begin{center}
\TEXTsymbol{\backslash}vspace\{\TEXTsymbol{\backslash}singlespace\}
\end{center}

\noindent The footnotes in this section have done these changes; you will
see the result in the compiled version.\footnote{%
%TCIMACRO{\TeXButton{singlespace}{\setlength{\baselineskip}{\singlespace}}}%
%BeginExpansion
\setlength{\baselineskip}{\singlespace}%
%EndExpansion
This is the second footnote on the page.}

\subsection{References}

There are two ways to handle references in LaTeX documents, typing them
directly or using BiBTeX. Examples of typing them directly are given below;
the gray TeX boxes immediately before the references are needed to make the
references be formatted correctly and have the references section show up in
the table of contents. No information on how to use BibTeX will be given
here; we only point out that to use BibTeX, you need to click on insert,
typeset object, bibliography, click on your BibTeX file and the appropriate
style. You should do this on the bottom line in main-sw.tex and not in
thesisbody.tex.

%TCIMACRO{\TeXButton{newpage}{\newpage\ }}%
%BeginExpansion
\newpage\ %
%EndExpansion
%TCIMACRO{%
%\TeXButton{addcontentsline REFERENCES}{\addcontentsline{toc}{section}{REFERENCES}}}%
%BeginExpansion
\addcontentsline{toc}{section}{REFERENCES}%
%EndExpansion
%TCIMACRO{%
%\TeXButton{setlength singlespace}{\setlength{\baselineskip}{\singlespace}}}%
%BeginExpansion
\setlength{\baselineskip}{\singlespace}%
%EndExpansion

\begin{thebibliography}{9}
\bibitem{Bez99} Guram Bezhanishvili, \emph{Varieties of monadic {H}eyting
algebras. {I}{I}. {D}uality theory}, Studia Logica \textbf{62} (1999),
no.~1, 21--48.

\bibitem{DeG45} J.~de~Groot, \emph{Topological classification of all closed
countable and continuous classification of all countable pointsets},
Indagationes Math. \textbf{7} (1945), 42--53.

\bibitem{Eng77} Ryszard Engelking, \emph{General topology}, PWN---Polish
Scientific Publishers, Warsaw, 1977.

\bibitem{Jon82} Peter~T. Johnstone, \emph{Stone spaces}, Cambridge
University Press, Cambridge, 1982.

\bibitem{Pri70} H.~A. Priestley, \emph{Representation of distributive
lattices by means of ordered stone spaces}, Bull. London Math. Soc. \textbf{2%
} (1970), 186--190.
\end{thebibliography}
